\documentclass[10pt]{article}
\usepackage[a4paper,margin=0.72in]{geometry}
\usepackage{amsmath,amssymb,amsthm,mathtools}
\usepackage{graphicx}
\usepackage[hidelinks]{hyperref}
\usepackage{xcolor}
\usepackage{microtype}
\setlength{\parindent}{0pt}
\setlength{\parskip}{4pt}
\newtheorem{theorem}{Theorem}
\newtheorem{lemma}[theorem]{Lemma}
\newtheorem{proposition}[theorem]{Proposition}
\newcommand{\TT}{\mathbb T}
\newcommand{\CC}{\mathbb C}
\newcommand{\W}{\mathcal W}
\newcommand{\JJ}{\mathcal J}
\newcommand{\supp}{\operatorname{supp}}

\title{Zero Jordan product determination for discrete and SIN group algebras}
\author{Candidate substantial partial solution for arXiv:1902.04846}
\date{27 August 2026}

\begin{document}
\maketitle

\section*{Result and relation to the source question}

Alaminos, Bre\v{s}ar, Extremera, and Villena prove that $L^1(G)$ is zero
Jordan product determined when $G$ is amenable, and ask whether this holds for
every locally compact group \cite[Question 3.7]{ABEV}:

\begin{center}
\includegraphics[width=.92\textwidth,trim=0 85 0 0,clip]{figures/open_problem_crop.png}
\end{center}

Recall that a locally compact group is a \emph{SIN group} if the identity has
a neighborhood basis invariant under conjugation.

\begin{theorem}\label{thm:main}
If $G$ is a locally compact SIN group, then $L^1(G)$ is zero Jordan product
determined.  In particular, $\ell^1(G)$ is zero Jordan product determined for
every discrete group $G$.
\end{theorem}

Thus there is no amenability, torsion, or finite-generation restriction in
the discrete case.  The result includes the torsion-free nonamenable free
groups.  It also includes nondiscrete nonamenable groups, for example
$F_2\times\mathbb R$.

\section*{The commutator criterion}

For a Banach algebra $A$, put $a\circ b=ab+ba$ and
\[
 \JJ(A)=\overline{\operatorname{span}}
 \{ab:a,b\in A,\ a\circ b=0\}.
\]
The source proves the following criterion: if $A$ is weakly amenable, has
property $\mathbb B$, and has a bounded approximate identity, then $A$ is
zero Jordan product determined if and only if
\[
 [A,A]\subseteq\JJ(A).                                      \tag{1}
\]
It also records that every group algebra $L^1(G)$ has all three hypotheses
\cite{ABEV}.  We will therefore prove (1).

The engine is an elementary synthesis statement for the Wiener algebra
\[
 \W=A(\TT)=\left\{f(z)=\sum_{n\in\mathbb Z}\widehat f(n)z^n:
       \|f\|_{\W}=\sum_n|\widehat f(n)|<\infty\right\}.
\]
Via Fourier coefficients,
$A(\TT^2)=\W\widehat\otimes_\pi\W$ isometrically.

\section*{A separated-support Wiener lemma}

\begin{lemma}\label{lem:wiener}
In $A(\TT^2)$ the function
\[
 D(z,w)=z-w
\]
belongs to the closed linear span of the tensors $g\otimes h$ with
$g,h\in\W$ and $\supp g\cap\supp h=\varnothing$.
\end{lemma}

\begin{proof}
We give the cutoff proof in full.  For small $\varepsilon>0$, choose a smooth
$2\pi$-periodic function $\chi_\varepsilon$ which is $1$ for
$|t|\leq\varepsilon$, is $0$ for $|t|\geq2\varepsilon$, and satisfies
$|\chi_\varepsilon'|\leq C/\varepsilon$.  Define
\[
 d_\varepsilon(e^{it})=(e^{it}-1)\chi_\varepsilon(t),\qquad
 q_\varepsilon(z)=(z-1)-d_\varepsilon(z).
\]
For every smooth $d$ on $\TT$, Cauchy--Schwarz on its Fourier coefficients
gives
\[
 \|d\|_{\W}\leq
 \left(\sum_{n\in\mathbb Z}(1+n^2)^{-1}\right)^{1/2}
 \|d\|_{H^1(\TT)}.                                         \tag{2}
\]
On an arc of length $O(\varepsilon)$ we have
$|e^{it}-1|=O(\varepsilon)$.  Hence
\[
 \|d_\varepsilon\|_2=O(\varepsilon^{3/2}),\qquad
 \|d_\varepsilon'\|_2=O(\varepsilon^{1/2}),
\]
and (2) implies $\|d_\varepsilon\|_{\W}\to0$.

Write $r=z\overline w$.  Since $D(z,w)=w(r-1)$, the smooth function
\[
 Q_\varepsilon(z,w)=wq_\varepsilon(z\overline w)
\]
vanishes on a neighborhood of the diagonal.  Moreover the change of Fourier
indices $(n\mapsto(n,1-n))$ gives
\[
 \|D-Q_\varepsilon\|_{A(\TT^2)}
 =\|d_\varepsilon\|_{\W}\longrightarrow0.                 \tag{3}
\]

It remains to decompose $Q_\varepsilon$.  Its compact support is covered by
finitely many product arcs $U_j\times V_j$ with
$\overline U_j\cap\overline V_j=\varnothing$.  Take a smooth partition of
unity on the support and write
$Q_\varepsilon=\sum_j Q_j$, where
$\supp Q_j\Subset U_j\times V_j$.  Choose smooth $\alpha_j,\beta_j$ supported
in $U_j,V_j$, respectively, and equal to $1$ on the corresponding coordinate
projections of $\supp Q_j$.  The double Fourier series
\[
 Q_j(z,w)=\sum_{m,n\in\mathbb Z}c^{(j)}_{mn}z^mw^n
\]
has absolutely summable coefficients.  Therefore, in projective norm,
\[
 Q_j=\alpha_j(z)\beta_j(w)Q_j(z,w)
 =\sum_{m,n}c^{(j)}_{mn}
       (\alpha_j(z)z^m)\otimes(\beta_j(w)w^n).              \tag{4}
\]
The two factors in every tensor in (4) have disjoint supports.  Equations
(3)--(4) prove the lemma.
\end{proof}

\section*{All discrete groups}

Let $G$ be any discrete group and $s\in G$.  There is a contractive unital
homomorphism
\[
 \pi_s:\W\longrightarrow\ell^1(G),\qquad
 \pi_s(f)=\sum_{n\in\mathbb Z}\widehat f(n)\delta_{s^n}.   \tag{5}
\]
This remains valid when $s$ has finite order: repeated point masses simply
combine, and absolute convergence proves both contractivity and
multiplicativity.

Fix $s,t\in G$ and define a bounded map
\[
 \Theta_{s,t}:A(\TT^2)\longrightarrow\ell^1(G),\qquad
 \Theta_{s,t}(g\otimes h)=\pi_s(g)\delta_t\pi_s(h).         \tag{6}
\]
If $\supp g\cap\supp h=\varnothing$, choose $u\in C^\infty(\TT)$ such that
$u=1$ on $\supp g$ and $u=-1$ on $\supp h$.  Set
\[
 a=\pi_s(u),\qquad b=\pi_s(g)\delta_t\pi_s(h).
\]
Since $ug=g$ and $hu=-h$, we have
\[
 ab=b,\qquad ba=-b.
\]
Thus $a\circ b=0$ and $b=ab\in\JJ(\ell^1(G))$.  Lemma
\ref{lem:wiener}, the continuity of (6), and
\[
 \Theta_{s,t}(D)=\delta_s\delta_t-\delta_t\delta_s
\]
show that every point-mass commutator belongs to $\JJ(\ell^1(G))$.

For arbitrary $p,q\in\ell^1(G)$,
\[
 [p,q]=\sum_{s,t\in G}p(s)q(t)[\delta_s,\delta_t]          \tag{7}
\]
converges absolutely in $\ell^1(G)$, since the sum of the norms is at most
$2\|p\|_1\|q\|_1$.  The space $\JJ(\ell^1(G))$ is closed and linear, so
(7) proves (1).  The source criterion completes the discrete case.

\section*{Transfer to locally compact SIN groups}

Let now $G$ be locally compact and SIN.  Shrinking relatively compact
conjugation-invariant neighborhoods $U$ of the identity give
\[
 e_U=m(U)^{-1}\mathbf1_U\in L^1(G).                         \tag{8}
\]
These form a contractive approximate identity.  They are central even in the
measure algebra $M(G)$.  Indeed, invariance of one such $U$ of finite positive
measure shows that inner automorphisms preserve Haar measure, so $G$ is
unimodular; conjugation invariance then gives
$\delta_s*e_U=e_U*\delta_s$ for every $s$, and integration gives commutation
with every measure.

Formula (5) now defines a contractive homomorphism $\pi_s:\W\to M(G)$.
Fix $s\in G$ and $k\in L^1(G)$, and put
\[
 \Theta_{s,k}(g\otimes h)=\pi_s(g)*k*\pi_s(h).             \tag{9}
\]
For a separated-support tensor, choose $u$ as above and set
\[
 a=\pi_s(u)\in M(G),\qquad b=\pi_s(g)*k*\pi_s(h)\in L^1(G).
\]
Again $a*b=b$ and $b*a=-b$.  Although $a$ need not lie in $L^1(G)$, define
$a_U=a*e_U\in L^1(G)$.  Centrality of $e_U$ gives the \emph{exact} identities
\[
 a_U*b=b*e_U,\qquad b*a_U=-b*e_U,
\]
so $a_U\circ b=0$ and $b*e_U=a_U*b\in\JJ(L^1(G))$.  Letting
$U$ shrink yields $b*e_U\to b$, hence $b\in\JJ(L^1(G))$.
Applying Lemma \ref{lem:wiener} to (9) now gives
\[
 [\delta_s,k]=\delta_s*k-k*\delta_s\in\JJ(L^1(G))         \tag{10}
\]
for all $s\in G$ and $k\in L^1(G)$.

If $f\in C_c(G)$, norm continuity of left and right translation makes the
following a Bochner integral in $L^1(G)$:
\[
 [f,k]=\int_G f(s)[\delta_s,k]\,ds.                        \tag{11}
\]
The integrand in (11) takes values in the closed linear space
$\JJ(L^1(G))$ by (10), so its Bochner integral does too.  Density of
$C_c(G)$ in $L^1(G)$ extends this to every $f\in L^1(G)$.  Thus (1) holds,
and the source criterion proves Theorem \ref{thm:main}.

\section*{Boundary of the method and novelty status}
\enlargethispage{3\baselineskip}

For a general locally compact group the cyclic calculus (5) still lives in
$M(G)$, but a bounded approximate identity need not be central.  Replacing
$a$ by $a*e_i$ then destroys the exact relation
$(a*e_i)\circ b=0$.  Approximate anticommutation is insufficient for the
definition of $\JJ(A)$.  We therefore do not claim a solution for arbitrary
locally compact groups.

A bounded novelty search on 27 August 2026 checked the run indexes and close
web/arXiv combinations of the source id with ``all discrete groups'', ``SIN
group'', ``central approximate identity'', ``Wiener algebra'', and
``diagonal synthesis''.  No explicit statement of either theorem above was
found.  This is provisional: an expert citation search is still required.

Human review should focus on projective-norm convergence in (4), the central
approximate identity in (8), and the multiplier-to-$L^1$ transfer.  These
steps are included explicitly to make the candidate independently auditable.

\begin{thebibliography}{9}
\bibitem{ABEV}
J. Alaminos, M. Bre\v{s}ar, J. Extremera, and A. R. Villena,
\emph{Zero Jordan product determined Banach algebras}, J. Aust. Math. Soc.
111 (2021), 145--158; arXiv:1902.04846.
\end{thebibliography}

\end{document}
